% GENERATED FILE. Edit canonical lessons, then run npm run generate:books.
% canonical-source-hash: fnv1a64:0fef11599fce6707
% canonical-lessons: ES-C407-medio, ES-C407-debajo, ES-C407-fuera, ES-C407-lado, ES-C407-location-recall-1, ES-C407-location-recall-2

\chapter{Middle, Below, Outside and Beside}
\label{ch:medio-debajo-fuera-lado}
\hlchaptermodality{407}

\begin{chapteropening}
\textbf{By the end of this chapter:} \emph{I can place something in the middle, underneath, outside or beside a reference point.}

Use four location frames without translating.
\end{chapteropening}

\section[en medio]{en medio — in the middle}

\label{lesson:ES-C407-medio}



\begin{quote}
Recall \textbf{la puerta}, the door. Now place it at the centre: \textbf{la puerta está en medio}.
\end{quote}

\begin{cousinweb}
\textbf{Medio} comes from Latin \textbf{medius}, “middle.” English kept the same root in \textbf{median}, \textbf{medium} and \textbf{intermediate}. Spanish uses \textbf{en medio} for a position between the edges or between two things.
\end{cousinweb}

\begin{grammarlens}[title={Grammar Lens: the frame en medio de}]
Say \textbf{en medio} when the reference is understood. Add \textbf{de} before the thing that frames the middle: \textbf{en medio de la puerta}, in the middle of the door.
\end{grammarlens}

\subsection*{Your turn}
Say these aloud:

\begin{itemize}
\raggedright
  \item \emph{en medio}
  \item \emph{en medio de la puerta}
  \item \emph{la puerta está en medio}
\end{itemize}

\begin{tcolorbox}[breakable,colback=teal!4,colframe=teal!35!black,title={Before you move on}]
In the middle? (\emph{En medio}.)

Source: \href{https://dle.rae.es/medio}{DLE: medio}.
\end{tcolorbox}
\section[debajo]{debajo — underneath}

\label{lesson:ES-C407-debajo}



\begin{quote}
Recall \textbf{la mesa}. Something lower than the tabletop can be \textbf{debajo de la mesa}.
\end{quote}

\begin{cousinweb}
The history is visible: \textbf{de + bajo}. \textbf{Bajo} is low; \textbf{debajo} names the lower position. The DLE defines it as a place below another higher one.
\end{cousinweb}

\begin{grammarlens}[title={Grammar Lens: debajo de}]
Before a reference object, use \textbf{debajo de}: \textbf{debajo de la mesa}. With no object stated, \textbf{está debajo} is complete.
\end{grammarlens}

\subsection*{Your turn}
Say these aloud:

\begin{itemize}
\raggedright
  \item \emph{debajo}
  \item \emph{debajo de la mesa}
  \item \emph{está debajo}
\end{itemize}

\begin{tcolorbox}[breakable,colback=teal!4,colframe=teal!35!black,title={Before you move on}]
Underneath? (\emph{Debajo}.) Under the table? (\emph{Debajo de la mesa}.)

Source: \href{https://dle.rae.es/debajo}{DLE: debajo}.
\end{tcolorbox}
\section[fuera]{fuera — outside}

\label{lesson:ES-C407-fuera}



\begin{quote}
Recall \textbf{la casa}. Inside it is one location; beyond it is \textbf{fuera}, outside.
\end{quote}

\begin{cousinweb}
\textbf{Fuera} comes from Latin \textbf{foras}, “outward.” The core picture is crossing the boundary and standing beyond it. The opening \textbf{fue-} is one syllable: say FWE-ra, not fu-E-ra.
\end{cousinweb}

\begin{grammarlens}[title={Grammar Lens: fuera and fuera de}]
\textbf{Está fuera} means it is outside. Add \textbf{de} for the boundary: \textbf{fuera de la casa}, outside the house.
\end{grammarlens}

\subsection*{Your turn}
Say these aloud:

\begin{itemize}
\raggedright
  \item \emph{fuera}
  \item \emph{está fuera}
  \item \emph{fuera de la casa}
\end{itemize}

\begin{tcolorbox}[breakable,colback=teal!4,colframe=teal!35!black,title={Before you move on}]
Outside? (\emph{Fuera}.)

Source: \href{https://dle.rae.es/fuera}{DLE: fuera}.
\end{tcolorbox}
\section[el lado]{el lado — the side}

\label{lesson:ES-C407-lado}



\begin{quote}
A door has a left edge and a right edge. Each edge belongs to \textbf{un lado}, a side.
\end{quote}

\begin{cousinweb}
\textbf{Lado} continues Latin \textbf{latus}, “side.” English \textbf{lateral} is the learned cousin: a lateral movement goes toward one side.
\end{cousinweb}

\begin{grammarlens}[title={Grammar Lens: al lado de}]
The useful location frame is \textbf{al lado de}, beside or next to: \textbf{al lado de la puerta}. The \textbf{al} is \textbf{a + el}, literally “at the side of.”
\end{grammarlens}

\subsection*{Your turn}
Say these aloud:

\begin{itemize}
\raggedright
  \item \emph{el lado}
  \item \emph{al lado}
  \item \emph{al lado de la puerta}
\end{itemize}

\begin{tcolorbox}[breakable,colback=teal!4,colframe=teal!35!black,title={Before you move on}]
The side? (\emph{El lado}.) Beside the door? (\emph{Al lado de la puerta}.)

Source: \href{https://dle.rae.es/lado}{DLE: lado}.
\end{tcolorbox}
\section[(review)]{location recall one — place four points}

\label{lesson:ES-C407-location-recall-1}



\begin{quote}
Give the four locations. (\emph{En medio, debajo, fuera, el lado}.)
\end{quote}

\subsection*{Your turn}
Say these aloud:

\begin{itemize}
\raggedright
  \item \emph{en medio de la puerta}
  \item \emph{debajo de la mesa}
  \item \emph{fuera de la casa}
  \item \emph{al lado de la puerta}
\end{itemize}

\begin{tcolorbox}[breakable,colback=teal!4,colframe=teal!35!black,title={Before you move on}]
Middle, below, outside, side? (\emph{Medio, debajo, fuera, lado}.)
\end{tcolorbox}
\section[(review)]{location recall two — complete the map}

\label{lesson:ES-C407-location-recall-2}



\begin{quote}
Complete each frame: \textbf{en \_\_\_ de}, \textbf{\_\_\_ de}, \textbf{\_\_\_ de}, \textbf{al \_\_\_ de}. (\emph{Medio, debajo, fuera, lado}.)
\end{quote}

\subsection*{Your turn}
Say these aloud:

\begin{itemize}
\raggedright
  \item \emph{en medio}
  \item \emph{debajo}
  \item \emph{fuera}
  \item \emph{al lado}
\end{itemize}

\begin{tcolorbox}[breakable,colback=teal!4,colframe=teal!35!black,title={Before you move on}]
Place an object in four different relations without translating. (\emph{En medio de, debajo de, fuera de, al lado de}.)
\end{tcolorbox}
